%============================================================ 5. 실험 설계
\section{실험 설계}
\label{sec:experiment}

\subsection{조건 행렬}

두 계층의 기여를 분리하려면 각 계층을 독립적으로 켜고 끌 수 있어야 한다. 따라서 배정 축과
기동 축을 각각 이분하여 $2\times2$ 조건 행렬을 구성한다(Table~\ref{tab:conditions}).

\begin{table*}[t]
\centering
\caption{$2\times2$ 조건 행렬. 조건 사이에 달라지는 것은 두 스위치뿐이며, 체크포인트·환경
         설정·난수 시드는 모두 동일하다.}
\label{tab:conditions}
\footnotesize
\begin{tabular}{llp{0.42\linewidth}}
\toprule
조건 & 배정 $\times$ 기동 & 무엇을 측정하는가 \\
\midrule
\texttt{heur\_heur} & 휴리스틱 $\times$ 휴리스틱 & 기준선(baseline) \\
\texttt{heur\_unet} & 휴리스틱 $\times$ U-Net    & \textbf{기동 계층의 순수 기여}
                                                   --- LLM을 쓰지 않는다 \\
\texttt{llm\_heur}  & LLM $\times$ 휴리스틱      & \textbf{전략 계층의 순수 기여} \\
\texttt{llm\_unet}  & LLM $\times$ U-Net         & 제안 체계(상호작용 포함) \\
\bottomrule
\end{tabular}
\end{table*}

휴리스틱 배정은 위협이 큰 무리부터 최근접 방어정을 탐욕적으로 1:1 대응시키되 직전 배정을
우대하는(sticky) 규칙이다. 휴리스틱 기동은 배정된 요격점에서 회랑을 가로지르는 방향으로
두 점을 잡는다. 중요한 점은 휴리스틱 기동이 정책과 \emph{동일한 유효 마스크와 거리 제약}을
거쳐 픽셀로 스냅된다는 것이다. 두 조건이 같은 행동공간을 쓰지 않으면 비교가 성립하지 않는다.

\texttt{heur\_unet} 조건은 특별한 위치를 갖는다. 이 조건은 언어모델을 전혀 호출하지 않으므로,
여기서 나온 기여는 외부 API에 의존하지 않고 재현 가능하다.

\subsection{지휘관 모델}

전략 계층에는 능력 수준이 다른 세 모델을 사용한다: Qwen2.5 7B, Qwen2.5 14B(둘 다 로컬
실행), 그리고 Gemini 3.5 Flash-Lite(클라우드 API). 로컬 모델을 포함한 이유는 해상 방어가
오프라인 운용을 요구할 수 있기 때문이며, 이 축은 단순한 성능 비교가 아니라
\emph{온보드 대 클라우드}라는 배치 선택지를 나타낸다.

\subsection{시드 대응표본 설계}
\label{sec:paired}

모든 조건이 \emph{동일한 시드 집합}(0--29)을 사용한다. 시드 $k$ 의 전장은 조건이 달라도
동일하므로, 조건 간 차이를 취하면 시나리오 난이도의 변동이 상쇄된다. 이는 독립표본 설계보다
훨씬 좁은 신뢰구간을 주어 적은 표본으로 결론에 도달하게 한다. 조건마다 다른 시드를 쓰면
이 설계의 전제가 무너져 통계가 무효가 되므로, 수집 파이프라인이 시드 집합을 조건에 걸쳐
강제한다.

총 수집량은 8조건 $\times$ 3포메이션 $\times$ 30시드 $=$ \textbf{720 에피소드}이다. 조건별
·포메이션별 수집이 균형을 이루는지 확인하였으며(모든 칸 30개), 불균형이 있으면 집계에서
공통 교집합만 사용하도록 하였다. 수집이 중간에 끊겨 어떤 조건이 쉬운 포메이션만 덮으면
그 조건이 부당하게 유리해지기 때문이다.

\subsection{재계획 주기}

지휘관 재계획은 매 결정(25\,step)마다 수행한다. 이는 실제 배치 설정과 동일한 주기이다.
재계획 주기를 늘리면 API 호출은 줄지만 배정 계획이 그만큼 낡는데, 휴리스틱 배정은 매 결정
새로 계산되므로 결과적으로 \emph{낡은 LLM 계획 대 최신 휴리스틱}을 비교하게 되어 LLM 조건에
불리한 편향이 생긴다. 이를 피하기 위해 주기를 배치 설정과 일치시키고, 원자료에 주기를
기록하여 집계 전 단일 값인지 확인하였다.

\subsection{통계 절차}
\label{sec:stats}

\textbf{신뢰구간.} 조건별 평균의 95\,\% 신뢰구간은 편향-가속 보정(BCa) 부트스트랩으로
구한다~\citep{efron1993bootstrap, diciccio1996bca}. 포획률은 $[0,1]$ 로 잘려 있고 분포가
한쪽으로 치우치므로, 단순 백분위 부트스트랩은 편향된 구간을 준다. BCa는 편향과 가속을
보정한다.

\textbf{대응 차이.} 기준선 대비 개선은 시드별 차이 $d_k = x_k - y_k$ 의 부트스트랩 분포로
평가한다. 신뢰구간이 0을 포함하면 그 조건이 기준선과 다르다고 말하지 않는다. 평균이 높아도
마찬가지이며, 본 논문은 그러한 경우를 그대로 보고한다.

\textbf{효과크기.} 대응표본 효과크기 $\dz = \bar{d} / s_d$ 를 함께 보고한다. 지표마다 단위가
다르므로 개선 폭의 절대값만으로는 지표 간 비교가 불가능하기 때문이다.
\S\ref{sec:results}에서 보이듯 이 비교가 본 논문의 핵심 관찰 중 하나로 이어진다.

\textbf{다중비교.} 한 조건에 대해 지표 10개를 동시에 검정하므로, 보정 없이 ``구간이
0을 배제''만 보면 우연히 유의해 보이는 지표가 섞인다. 본 논문은 두 가지를 함께 보고한다.
보정하지 않은 95\,\% 구간과, 지표 10개에 대한 Bonferroni 보정 구간(99.5\,\%)이다
(Table~\ref{tab:effectsize}). 조건 수까지 곱하지 않은 것은 조건별 비교가 각각 독립된
질문이기 때문이며, 이 선택을 여기에 밝혀 숨은 자유도를 남기지 않는다. 어느 결론이 보정을
견디고 어느 것이 못 견디는지는 \S\ref{sec:res_effectsize}에서 그대로 보고한다.

\textbf{상호작용 분해.} 두 계층이 서로를 돕는지를 판정하기 위해 $2\times2$ 상호작용을
대응표본으로 분해한다.
\begin{align}
\mathrm{Int} =\;& (\texttt{llm\_unet} - \texttt{llm\_heur}) \nonumber\\
                &- (\texttt{heur\_unet} - \texttt{heur\_heur})
\label{eq:interaction}
\end{align}
$\mathrm{Int} > 0$ 이면 두 계층이 서로를 돕고(상승), $\approx 0$ 이면 두 기여가 단순
합산(가산)이며, $< 0$ 이면 서로를 잠식(상쇄)한다. 본 논문은 이 값을 모든 지표에 대해
계산하여 보고한다.

\subsection{지휘관 계층 계측}

포획률은 체계 전체의 지표이므로 언어모델 계층이 잘했는지를 가린다. 따라서 시뮬레이션 결과와
\emph{독립적으로}, 계획이 산출된 시점에 확인 가능한 것만 따로 계측한다: 응답 지연,
폴백률(언어모델 실패로 휴리스틱으로 대체된 비율), 커버리지(활성 무리 중 배정된 비율),
churn(재계획 간 담당이 바뀐 방어정 비율), 교차 배정(두 방어정의 요격 경로가 교차하는 비율).

폴백률은 특히 중요하다. 이 값이 0이 아니면 해당 ``LLM 조건''의 일부가 사실은 휴리스틱
출력이며, 그 상태의 수치는 언어모델의 성능이 아니다. 본 수집에서 폴백률은 Gemini 0.000,
Qwen2.5 14B 0.001, Qwen2.5 7B 0.007로 모두 무시할 수 있는 수준이었다.
